% GraphHunter — Libro «Diseño completo»
% Pedagogical design doc: complete walkthrough of GraphHunter's
% architecture, algorithms, and forward-looking stages.
%
% Accent color: morado oscuro (Diseño) for visual distinction
% from the other four books.

\usepackage{fontspec}
\setmainfont{Latin Modern Roman}
\setsansfont{Latin Modern Sans}
\setmonofont{Latin Modern Mono}
\usepackage[spanish,es-tabla,es-lcroman,es-noshorthands]{babel}
\usepackage{csquotes}

\usepackage[a4paper,margin=2.4cm]{geometry}

\usepackage{amsmath, amssymb, amsthm}
\usepackage{mathtools}

\usepackage[ruled,linesnumbered,algochapter]{algorithm2e}
\SetKwComment{Comment}{$\triangleright$ }{}
\SetKwFor{For}{for}{do}{end}
\SetKwFor{While}{while}{do}{end}
\SetKwIF{If}{ElseIf}{Else}{if}{then}{else if}{else}{end}
\SetKwRepeat{Repeat}{repeat}{until}

\usepackage{listings}
\usepackage{xcolor}
\definecolor{accent}{HTML}{4B2E83}       % morado oscuro (Diseño)
\definecolor{accentlight}{HTML}{EAE0F4}
\definecolor{codecomment}{rgb}{0.35,0.45,0.45}
\definecolor{codekw}{rgb}{0.55,0.10,0.55}
\definecolor{codestr}{rgb}{0.10,0.45,0.10}
\lstdefinelanguage{Rust}{
    morekeywords={fn,let,mut,pub,struct,enum,impl,trait,where,if,else,for,while,
                  loop,match,return,use,mod,const,static,self,Self,as,in,ref,move,
                  async,await,dyn,unsafe,crate,super,extern,type,break,continue},
    morecomment=[l]{//},
    morecomment=[s]{/*}{*/},
    morestring=[b]",
    morestring=[b]',
}
\lstset{
    basicstyle=\ttfamily\small,
    keywordstyle=\color{codekw}\bfseries,
    commentstyle=\color{codecomment}\itshape,
    stringstyle=\color{codestr},
    showstringspaces=false,
    frame=single,
    numbers=none,
    captionpos=b,
    breaklines=true,
    tabsize=4,
    % Listings comen Unicode salvo que les digas literate→glyph.
    % Mapeamos los símbolos griegos / matemáticos que aparecen en
    % comentarios de Rust dentro de los listings.
    extendedchars=true,
    literate=%
        {α}{{$\alpha$}}1
        {β}{{$\beta$}}1
        {ρ}{{$\rho$}}1
        {Σ}{{$\Sigma$}}1
        {≤}{{$\leq$}}1
        {≥}{{$\geq$}}1
        {→}{{$\rightarrow$}}1
        {↔}{{$\leftrightarrow$}}1,
}

\usepackage{hyperref}
\hypersetup{
    colorlinks=true,
    linkcolor=accent,
    citecolor=black,
    urlcolor=accent,
    pdftitle={GraphHunter — Diseño completo},
    pdfauthor={Lucas Sotomayor / Base4Security},
    bookmarks=true,
    bookmarksopen=true,
    bookmarksnumbered=true,
}
\usepackage{bookmark}

\usepackage{booktabs}
\usepackage{longtable}
\usepackage{float}
\usepackage{tikz}
\usetikzlibrary{arrows.meta, positioning, shapes.geometric, fit, backgrounds, calc, shapes.misc, decorations.pathreplacing}

% Box style for "intuición" callouts — used to give a plain-language
% summary before each formal definition.
\usepackage[most]{tcolorbox}
\newtcolorbox{intuicion}{
    colback=accentlight,
    colframe=accent,
    arc=2mm,
    boxrule=0.6pt,
    left=4pt, right=4pt, top=4pt, bottom=4pt,
    fonttitle=\bfseries,
    title=Intuición
}

\theoremstyle{plain}
\newtheorem{theorem}{Teorema}[chapter]
\newtheorem{lemma}[theorem]{Lema}

\theoremstyle{definition}
\newtheorem{definition}[theorem]{Definición}

\theoremstyle{remark}
\newtheorem*{invariant}{Invariante}
\newtheorem*{nota}{Nota}

\newcommand{\filepath}[1]{\texttt{\small\detokenize{#1}}}
\newcommand{\code}[1]{\texttt{\detokenize{#1}}}

\newenvironment{enelcodigo}{%
    \par\vspace{0.4em}%
    \noindent\textbf{En el código.}\nopagebreak\par\nopagebreak\vspace{0.2em}%
}{\par\vspace{0.4em}}

\usepackage{titlesec}
\titleformat{\chapter}[display]
  {\normalfont\huge\bfseries\color{accent}}
  {\chaptertitlename~\thechapter}
  {20pt}
  {\Huge}

\newcommand{\TODO}[1]{\textbf{\textcolor{red}{[TODO: #1]}}}
\newcommand{\passthrough}[1]{#1}
